\section{\textbf{From General Intelligence to Superintelligence}}
\label{sec:agi_to_si}

Section~\ref{sec:introduction} proposed that superintelligence begins when a
generally intelligent machine originates validated knowledge beyond the
collective human frontier. This section develops that distinction more
precisely. The transition from AGI to superintelligence is not defined by a
larger model, a longer context window, access to more tools, or improved
performance within an existing task. It is defined by a change in the machine's
epistemic role: from a system that operates across established knowledge and
human-specified problems to a system that can expose, formulate, and resolve
problems that humanity had not recognized.

\subsection{\textbf{The Functional Boundary of General Intelligence}}

General intelligence concerns the organization of agency. An AGI must transfer
knowledge across domains, recompose existing capabilities for unfamiliar
problems, maintain goals over time, revise plans after failure, and coordinate
perception, memory, reasoning, metacognition, and action
\cite{zohourian2026agi}. Formal definitions of intelligence likewise emphasize
the capacity to achieve goals across a wide range of environments
\cite{legg2007universal}, while broader AGI accounts distinguish general
adaptability from competence on isolated tasks \cite{goertzel2014agi}.

Let $\mathcal{E}$ denote environments, $\mathcal{T}$ tasks, and $\mathcal{M}$
modalities. For agent $x$, let $T_x$, $R_x$, and $A_x$ denote transfer,
recomposition, and adaptation. The functional scope of general agency can be
represented as

\begin{equation}
\begin{aligned}
\mathcal{G}(x)
&= \Phi_x(T_x,R_x,A_x),\\
\Phi_x:
\mathcal{E}\times\mathcal{T}\times\mathcal{M}
&\longrightarrow
\text{coherent goal-directed behavior},
\end{aligned}
\label{eq:general_agency}
\end{equation}

where $\Phi_x$ denotes the persistent organization through which these
capacities inform and constrain one another. No single component is sufficient.
Transfer without adaptation may fail when prior assumptions break;
recomposition without coherent goals may produce locally useful but globally
inconsistent actions; and broad task exposure may reflect memorized coverage
rather than general agency.

Generality also cannot be inferred from the mere presence of many components.
A framework may contain language, vision, memory, search, code execution, and
robotic interfaces while relying on a human to select tools, reconcile their
outputs, detect failure, and preserve the objective. Framework capacity,
realized agency, and observed performance must therefore remain distinct:

\begin{equation}
\begin{aligned}
\operatorname{Cap}(F)
&\not\equiv \operatorname{Realize}(x,F),\\
\operatorname{Realize}(x,F)
&\not\equiv \operatorname{Perf}(x,\tau),
\end{aligned}
\label{eq:capacity_realization}
\end{equation}

where $F$ is a computational framework and $\tau$ is an evaluated task. An AGI
realizes coordinated general agency; it need not be omniscient, infallible, or
superior on every task. Its defining property is the ability to reorganize its
available capacities in pursuit of objectives across sufficiently varied and
unfamiliar conditions.

\subsection{\textbf{The Epistemic Boundary of General Intelligence}}

Nothing in the preceding definition requires AGI to create knowledge beyond
humanity. A generally intelligent machine could possess extensive human
knowledge, apply it across disciplines, generate plans, conduct analyses, and
solve previously unsolved but already formulated problems. Such a system could
be transformative while remaining epistemically bounded by the concepts,
questions, and objectives inherited from human inquiry.
Reviews of AI-enabled scientific discovery already distinguish multiple roles
for machines across hypothesis generation, modelling, experimental design, and
analysis \cite{wang2023scientific}. That breadth of participation motivates a
more precise question of attribution: whether the machine merely operates
within the human research agenda or originates an epistemic object beyond it.

Let $\mathcal{K}_{H,t}$ be collective justified human knowledge at time $t$,
and let $\mathcal{Q}_{H,t}$ be the set of questions, hypotheses, and research
targets recognized by humans. An AGI operating within the human epistemic
boundary may be represented by

\begin{equation}
\begin{aligned}
x:
(\mathcal{K}_{H,t},\mathcal{Q}_{H,t},\mathcal{O}_{t})
&\longrightarrow
(\mathcal{Y}_{t},\mathcal{A}_{t}),\\
\mathcal{Q}_{x,t}
&\subseteq \mathcal{Q}_{H,t},
\end{aligned}
\label{eq:agi_epistemic_boundary}
\end{equation}

where $\mathcal{O}_{t}$ denotes available observations,
$\mathcal{Y}_{t}$ generated conclusions or predictions, and
$\mathcal{A}_{t}$ actions. The subset relation expresses the limiting case in
which the machine's research agenda remains contained within questions humans
have already recognized. The machine may produce an answer absent from
$\mathcal{K}_{H,t}$ and thereby resolve a known unknown, but it does not yet
originate the object of inquiry itself.

This distinction is visible in current AI-assisted science. Machine learning
has helped recover symbolic laws from experimental data
\cite{schmidt2009laws,udrescu2020aifeynman}, guided
mathematicians toward conjectures and theorems \cite{davies2021mathematics},
predicted protein structures \cite{jumper2021alphafold}, and proposed large
numbers of candidate stable materials \cite{merchant2023materials}. Language
models connected to tools and laboratory automation can design, plan, and
execute chemistry experiments \cite{boiko2023autonomous}, while autonomous
laboratories can perform closed-loop materials synthesis
\cite{szymanski2023autonomous}. Earlier robot-scientist systems likewise
generated and experimentally tested functional-genomics hypotheses
\cite{king2004robot} and demonstrated end-to-end automation of iterative
scientific investigation \cite{king2009automation}. These are consequential advances at the
scientific frontier. They also begin from domains, objectives, datasets,
experimental systems, and validation conditions established by human
researchers. Their results demonstrate increasingly powerful scientific
instruments and agents, not by themselves the invention of an unknown unknown.

\subsection{\textbf{Autonomous Problem Formation}}

The first additional requirement for superintelligence is autonomous problem
formation. The machine must be capable not only of answering a question but of
identifying that a consequential question exists. This may occur when the
system detects an unexplained regularity, recognizes a contradiction across
disciplines, constructs a new conceptual variable, identifies an assumption
shared by existing theories, or notices that available observations support a
previously unrepresented phenomenon. Work on automated identification of
latent state variables from high-dimensional observations illustrates progress
toward machines constructing representations that were not specified directly
by researchers \cite{chen2022variables}. Such representation discovery is an
important precursor, although it does not alone establish autonomous problem
formation in the stronger sense proposed here.

Let $q_{x}$ be a machine-originated research question. Autonomous problem
formation requires

\begin{equation}
\begin{aligned}
q_x \notin \mathcal{Q}_{H,t},
\qquad
q_x \notin \mathcal{K}_{H,t},\\
\operatorname{Ground}(q_x,\mathcal{K}_{H,t},\mathcal{O}_{t})=1,
\qquad
\operatorname{Consequential}(q_x)=1.
\end{aligned}
\label{eq:problem_formation}
\end{equation}

The grounding condition prevents arbitrary question generation from being
mistaken for intelligence. A language model can generate indefinitely many
unusual questions through stochastic variation. Most will be trivial,
ill-posed, physically impossible, semantically incoherent, or disconnected
from evidence. A superintelligent problem formulation must arise from a
defensible relation among observations, established knowledge, unresolved
tensions, or formally represented possibilities. The consequentiality
condition further requires that resolving the question would materially alter
understanding, explanation, prediction, or technological capability.

Autonomous problem formation is therefore more demanding than novelty. It
requires epistemic judgment: the system must distinguish an unasked question
that matters from an unasked question that is merely possible. This capacity
is the point at which an unknown unknown becomes an identified object of
inquiry.

\subsection{\textbf{Causal Origination of Knowledge}}

Scientific results are commonly produced by distributed systems containing
researchers, datasets, instruments, software, models, and institutions.
Requiring a superintelligent machine to work without human participation would
therefore impose a standard not applied to human scientists. Nevertheless,
attribution cannot be granted merely because AI appears somewhere in the
workflow. The machine must make the causally essential epistemic contribution.

Let $W$ be a scientific workflow and let $\kappa$ be its resulting knowledge
claim. We define machine origination counterfactually:

\begin{equation}
\begin{aligned}
\operatorname{Originate}_{x}(\kappa\mid W)=1
\Longleftrightarrow{}&
\operatorname{Produce}(W,\kappa)=1\\
&\land
\operatorname{Produce}(W\setminus x,\kappa)=0\\
&\land
\operatorname{Specified}_{H}(\kappa)=0.
\end{aligned}
\label{eq:causal_origination}
\end{equation}

The first condition requires that the workflow produce the contribution. The
second requires that removing the machine's relevant epistemic operation would
prevent that contribution from being produced through the same workflow. The
third excludes cases in which humans fully specified the conclusion and used
the machine only to calculate, retrieve, format, or experimentally execute it.
This is a conceptual causal test rather than a demand for literal repetition of
every discovery process.

Machine origination can occur through knowledge fusion or knowledge discovery.
In fusion, the essential contribution may be recognizing that previously
separated theories, data, or disciplines instantiate a common structure. In
discovery, it may be detecting a new entity, mechanism, regularity, or law.
Neither pathway is sufficient until its output is converted into an epistemic
object that can be understood, assessed, and incorporated into knowledge.
Prior systems for symbolic-law recovery and automated hypothesis testing
demonstrate that portions of this cycle can be computationally realized
\cite{schmidt2009laws,king2004robot}; the unresolved issue is whether the
machine can originate the consequential object of inquiry rather than search
within a human-defined hypothesis space.

\subsection{\textbf{The Transition Criterion}}

The transition from AGI to superintelligence requires a complete epistemic
cycle. The machine must form or expose the previously unrecognized problem,
originate a candidate explanation or solution, connect that contribution to
existing knowledge, and enable domain-appropriate validation. Let
$\kappa_x$ denote the resulting machine-originated contribution. Then

\begin{equation}
\begin{aligned}
\operatorname{KI}_{x}(\kappa_x)
\Longleftrightarrow{}&
\operatorname{FormProblem}_{x}(q_x)\\
&\land \operatorname{Originate}_{x}(\kappa_x\mid W)\\
&\land \operatorname{Explain}_{x}
(\kappa_x,\mathcal{K}_{H,t})\\
&\land \operatorname{Validate}(\kappa_x)=1\\
&\land \operatorname{Consequential}(\kappa_x)=1.
\end{aligned}
\label{eq:knowledge_invention_cycle}
\end{equation}

Knowledge invention is not established at the moment a machine produces a
novel output. Before validation, the output is a candidate hypothesis, theory,
method, or invention. Validation may require controlled experiment, formal
proof, reproducible computation, engineering realization, convergent evidence,
or another procedure accepted by the relevant discipline. Independent human or
machine verification may participate in this process, but the validation
standard must not be defined solely by the originating system. This insistence
on domain-appropriate validation follows the broader AI-for-science literature,
in which prediction, hypothesis generation, experimentation, and analysis form
an interdependent discovery process rather than interchangeable evidence of
knowledge \cite{wang2023scientific}.

The relationship between AGI and superintelligence can consequently be stated
as

\begin{equation}
\begin{aligned}
\operatorname{SI}(x)
\Longleftrightarrow{}&
\operatorname{AGI}(x)\\
&\land \exists q_x,\kappa_x:
q_x\notin\mathcal{Q}_{H,t}\\
&\land \operatorname{KI}_{x}(\kappa_x)\\
&\land
\mathcal{K}_{H,t+1}
=\mathcal{K}_{H,t}\cup\{\kappa_x\}.
\end{aligned}
\label{eq:si_transition}
\end{equation}

Equation~\eqref{eq:si_transition} preserves the established meaning of
superintelligence as capability beyond humans \cite{bostrom2014superintelligence}
while replacing an unspecified performance comparison with an epistemic
criterion. The machine goes beyond humans not merely by doing a known task
better, but by making a validated addition to collective knowledge whose object
of inquiry humanity had not recognized.

This criterion also clarifies why the present framework does not define a
separate category of domain superintelligence. A specialized scientific model
may be indispensable to a discovery, but specialization does not establish
general intelligence. Under the terminology developed here,
superintelligence is a mode of general intelligence. Specialized systems are
scientific instruments or components of a discovery system unless they also
realize the general agency and machine-originated knowledge invention required
by Equation~\eqref{eq:si_transition}.

Knowledge fusion and knowledge discovery now require separate analysis because
they describe the principal routes by which a machine may cross the human
epistemic boundary. The following sections examine how each route can produce
knowledge invention, how their outputs differ from retrieval and recombination,
and what evidence would justify attributing the resulting knowledge to the
machine.
